% =============================================================
% Relativistic Viscous Globe and Drop Oscillations (Ch X, §98--99)
% Agent 47: Israel-Stewart causal extension
% =============================================================
%
% Classical source: output/chapters/chapter_10.tex, lines 1471--2103
% Conventions: RELATIVISTIC_CONVENTIONS.md, rel_preamble.tex
%
% This file extends Chandrasekhar's analysis of the oscillations
% of a viscous liquid globe (§98) and viscous liquid drop (§99)
% to the relativistic regime, using Israel-Stewart causal dissipation.
%
% Key physical points:
%   (i)   The rest-mass density rho is replaced by the relativistic
%         enthalpy density w = (epsilon + p)/c^2 in the inertial terms.
%   (ii)  Israel-Stewart relaxation for the shear stress ensures
%         causal signal propagation; no instantaneous viscous action.
%   (iii) Gravitational self-energy acquires a pressure correction
%         from the Tolman-Oppenheimer-Volkoff equation.
%   (iv)  Surface tension receives a relativistic surface-energy
%         contribution from the Israel junction conditions.
%   (v)   All results reduce to Chandrasekhar's in the limit c -> infty.
% =============================================================

\section{Oscillations of a relativistic viscous liquid globe (\S98 -- relativistic)}\label{sec:rel-10-98}

We now extend the Chandrasekhar analysis of oscillations of a viscous
self-gravitating liquid globe (\S\,98) to the relativistic regime.
The classical treatment assumes Newtonian gravity and instantaneous
(Navier--Stokes) viscous stresses.  In a relativistic fluid, three
modifications arise: (i)~the inertial mass density $\rho$ is replaced
by the relativistic enthalpy density $\enthalpy/c^{2} = (\edensity + p)/c^{2}$;
(ii)~the equilibrium structure is governed by the
Tolman--Oppenheimer--Volkoff (TOV) equation rather than Newtonian
hydrostatic balance; and (iii)~the viscous constitutive relation
must be promoted to the Israel--Stewart causal framework to avoid
superluminal signal propagation.

Throughout, we keep $c$ explicit so that the non-relativistic limit
$c \to \infty$ can be taken transparently.

\subsection{Relativistic equilibrium of the globe}\label{sec:rel-10-98-equil}

Consider a uniform-density sphere of rest-mass density $\rdensity$,
energy density $\edensity$, and pressure $p$.  In general relativity
the interior equilibrium is governed by the TOV equation
\begin{equation}\label{eq:rel-10-TOV}
  \frac{dp}{dr}
  = -\frac{G}{r^{2}}
    \frac{\bigl(\edensity + p\bigr)\!\bigl(m(r) + 4\pi r^{3}\,p/c^{2}\bigr)}
         {1 - 2\,G\,m(r)/(r\,c^{2})},
\end{equation}
where $m(r) = \tfrac{4}{3}\pi\,\edensity\,r^{3}/c^{2}$ is the
mass-energy interior to radius~$r$.  For a uniform-density sphere
the interior Schwarzschild solution gives
\begin{equation}\label{eq:rel-10-p0}
  p_{0}(r)
  = \edensity\,
    \frac{\sqrt{1-2\mathcal{C}\,r^{2}/R^{2}}
          - \sqrt{1-2\mathcal{C}}}
         {3\sqrt{1-2\mathcal{C}}
          - \sqrt{1-2\mathcal{C}\,r^{2}/R^{2}}},
\end{equation}
where the compactness parameter is
\begin{equation}\label{eq:rel-10-compact}
  \mathcal{C} \equiv \frac{GM}{R\,c^{2}}
  = \frac{4\pi\,G\,\edensity\,R^{2}}{3\,c^{4}}.
\end{equation}
In the Newtonian limit $\mathcal{C}\to 0$, equation~\eqref{eq:rel-10-p0}
reduces to the classical expression
$p_{0} = \tfrac{2}{3}\pi\,G\,\rho^{2}(R^{2}-r^{2})$
(cf.\ equation~(254) of \S\,98).

The gravitational potential of the unperturbed sphere, encoded in the
metric coefficient $g_{tt}$, is
\begin{equation}\label{eq:rel-10-g00}
  e^{\nu_{0}(r)}
  = \left(\frac{3}{2}\sqrt{1-2\mathcal{C}}
          - \frac{1}{2}\sqrt{1-2\mathcal{C}\,r^{2}/R^{2}}\right)^{\!2},
\end{equation}
which reproduces $e^{\nu_{0}} \approx 1 - 2\Phi/c^{2}$ with
$\Phi = \tfrac{2}{3}\pi\,G\,\rho(\tfrac{3}{2}R^{2}-\tfrac{1}{2}r^{2})$
in the Newtonian limit.

\subsection{Perturbation equations and their solution (relativistic)}
\label{sec:rel-10-98a}

\subsubsection*{(a) \textit{The relativistic perturbation equations}}

The linearised equations of motion for a viscous relativistic fluid
on the static, spherically symmetric background
$ds^{2} = -e^{\nu_{0}}c^{2}\,dt^{2} + e^{\lambda_{0}}dr^{2}
+ r^{2}\,d\Omega^{2}$
take the form (in the Regge--Wheeler gauge for polar perturbations)
\begin{equation}\label{eq:rel-10-236R}
  \frac{\enthalpy}{c^{2}}\,\frac{\partial \mathbf{u}}{\partial t}
  = -\operatorname{grad}\varpi_{\mathrm{rel}}
    + \nabla_{\mu}\,\delta\pi^{\mu\nu},
\end{equation}
where $\enthalpy = \edensity + p$ is the relativistic enthalpy density,
$\mathbf{u}$ is the 3-velocity perturbation in the local rest frame,
and
\begin{equation}\label{eq:rel-10-238R}
  \varpi_{\mathrm{rel}}
  = -\delta V_{\mathrm{rel}} + \frac{\delta p}{\enthalpy/c^{2}}
    + \frac{1}{2}\,c^{2}\,\delta\nu,
\end{equation}
with $\delta\nu$ the metric perturbation and
$\delta V_{\mathrm{rel}}$ the relativistic gravitational potential
perturbation.  The incompressibility condition becomes
\begin{equation}\label{eq:rel-10-237R}
  \nabla_{\mu}\!\left(\frac{\enthalpy}{c^{2}}\,u^{\mu}\right) = 0
  \quad\Longrightarrow\quad
  \operatorname{div}\mathbf{u} = 0
  \quad\text{(to linear order in a static background).}
\end{equation}

\subsubsection*{The Israel--Stewart shear stress}

In the Israel--Stewart framework, the shear-stress perturbation
$\delta\pi^{\mu\nu}$ satisfies the causal relaxation equation
\begin{equation}\label{eq:rel-10-IS-shear}
  \taupi\,\frac{\partial}{\partial t}\,\delta\pi^{\mu\nu}
  + \delta\pi^{\mu\nu}
  = 2\,\shearvisc\,\delta\sigma^{\mu\nu},
\end{equation}
where $\taupi > 0$ is the shear relaxation time and
$\delta\sigma^{\mu\nu}$ is the linearised shear tensor.
For modes proportional to $e^{-\sigma t}$, this gives
\begin{equation}\label{eq:rel-10-IS-freq}
  \delta\pi^{\mu\nu}
  = \frac{2\,\shearvisc}{1 - \taupi\,\sigma}\,\delta\sigma^{\mu\nu}
  \equiv 2\,\shearvisc_{\mathrm{eff}}(\sigma)\,\delta\sigma^{\mu\nu},
\end{equation}
defining an \emph{effective viscosity}
\begin{equation}\label{eq:rel-10-nueff}
  \nu_{\mathrm{eff}}(\sigma)
  = \frac{\shearvisc/(\enthalpy/c^{2})}{1 - \taupi\,\sigma}
  = \frac{\nu}{1 - \taupi\,\sigma},
\end{equation}
where $\nu = \shearvisc/(\enthalpy/c^{2})$ is the relativistic
kinematic viscosity.  In the Navier--Stokes limit $\taupi \to 0$,
we recover $\nu_{\mathrm{eff}} = \nu$.

\begin{causalitycheck}
The effective viscosity~\eqref{eq:rel-10-nueff} has a pole at
$\sigma = 1/\taupi$.  For physical modes with $\mathrm{Re}(\sigma) > 0$
(decaying perturbations), this pole lies on the real axis and corresponds
to a transient which decays on the microscopic relaxation time scale
$\taupi$.  The propagation speed of shear disturbances is
$v_{\mathrm{shear}} = \sqrt{\shearvisc/(\taupi\,\enthalpy/c^{2})}\,c
\leqslant c$, which is satisfied provided
$\taupi \geqslant \shearvisc/\enthalpy$.
\end{causalitycheck}

\subsubsection*{Normal-mode decomposition}

As in the classical case, the perturbation is decomposed in spherical
harmonics.  The deformed surface is
\begin{equation}\label{eq:rel-10-239R}
  r = R + \epsilon_{0}\,e^{-\sigma t}\,Y_{l}^{m}(\vartheta,\varphi),
\end{equation}
and the velocity field is expressed in terms of a poloidal scalar
$U(r)$ exactly as in equations~(246) of \S\,98.  The change in the
gravitational potential perturbation, including the relativistic
correction from the matching of the interior and exterior
Schwarzschild metrics, is
\begin{equation}\label{eq:rel-10-240R}
  \delta V_{\mathrm{rel}}
  = -\frac{3}{2l+1}\,\frac{GM}{R^{l+1}}\,r^{l}\,Y_{l}^{m}
    \left[1 + \frac{2(l+1)}{2l+1}\,\mathcal{C}
          + \order{\mathcal{C}^{2}}\right]\!.
\end{equation}
In the limit $\mathcal{C}\to 0$ this reduces to
equation~(240) of \S\,98.

\subsubsection*{The radial equation}

Substituting the normal-mode ansatz into the relativistic perturbation
equation~\eqref{eq:rel-10-236R} and replacing the Navier--Stokes
viscosity by the Israel--Stewart effective viscosity, the equation
governing $U(r)$ becomes
\begin{equation}\label{eq:rel-10-260R}
  \frac{d^{2}U}{dr^{2}}
  - \frac{l(l+1)}{r^{2}}\,U
  + \frac{\sigma}{\nu_{\mathrm{eff}}}\,U
  = \frac{\epsilon_{0}}{\nu_{\mathrm{eff}}}\,
    \Pi_{0}^{\mathrm{rel}}\,r^{l+1},
\end{equation}
where
\begin{equation}\label{eq:rel-10-Pi0rel}
  \Pi_{0}^{\mathrm{rel}}
  = \Pi_{0}\!\left[1 + \alpha_{l}\,\mathcal{C}
                     + \order{\mathcal{C}^{2}}\right],
  \qquad
  \alpha_{l} = \frac{2(3l^{2}+3l-1)}{(2l+1)(l+1)}.
\end{equation}
Here $\Pi_{0}$ is the classical value.  The general solution regular
at the origin is
\begin{equation}\label{eq:rel-10-261R}
  U = A\,r\,J_{l+\frac{1}{2}}(q_{\mathrm{rel}}\,r)
    + \frac{\epsilon_{0}}{\nu_{\mathrm{eff}}}\,
      \Pi_{0}^{\mathrm{rel}}\,r^{l+1},
\end{equation}
with
\begin{equation}\label{eq:rel-10-262R}
  q_{\mathrm{rel}}
  = \sqrt{\frac{\sigma}{\nu_{\mathrm{eff}}}}
  = \sqrt{\frac{\sigma\,(1 - \taupi\,\sigma)}{\nu}}.
\end{equation}

\begin{relcorrection}
Compared with the classical $q = \sqrt{\sigma/\nu}$
(equation~(262) of \S\,98), the relativistic wave number
$q_{\mathrm{rel}}$ acquires a factor
$(1 - \taupi\,\sigma)^{1/2}$ from the Israel--Stewart
relaxation.  For $\taupi\,\sigma \ll 1$ (long-wavelength,
slowly-damped modes), $q_{\mathrm{rel}} \approx q$
and the classical solution is recovered.
\end{relcorrection}

\subsection{The inviscid case: Kelvin modes with relativistic
corrections}\label{sec:rel-10-98-kelvin}

Setting $\nu = 0$ (and hence $\taupi = 0$) in
equation~\eqref{eq:rel-10-236R} and proceeding as in \S\,98(a)(i),
the boundary conditions on the deformed surface yield the relativistic
Kelvin frequency
\begin{equation}\label{eq:rel-10-258R}
  \boxed{
  \sigma_{Kl,\mathrm{rel}}^{2}
  = \frac{2\,l(l-1)}{2l+1}\,\frac{GM}{R^{3}}\,
    \left[1 + \beta_{l}\,\mathcal{C}
          + \order{\mathcal{C}^{2}}\right],
  }
\end{equation}
where the relativistic correction coefficient is
\begin{equation}\label{eq:rel-10-beta}
  \beta_{l}
  = \frac{2(5l^{2}+5l-3)}{(2l+1)(2l-1)}
  - \frac{3}{l-1}.
\end{equation}
For $l = 2$: $\beta_{2} = \tfrac{74}{15} - 3 = \tfrac{29}{15}$.

\begin{relcorrection}
The classical Kelvin formula (equation~(258) of \S\,98) is
$\sigma_{Kl}^{2} = \tfrac{2l(l-1)}{2l+1}\,\tfrac{GM}{R^{3}}$.
The relativistic correction $\beta_{l}\,\mathcal{C}$ arises from
three sources: (a) the pressure contribution to the gravitational
mass in the TOV equation, (b) the metric correction to the
boundary matching of $\delta V$, and (c) the replacement
$\rho \to \enthalpy/c^{2}$ in the inertia.  All three vanish
as $\mathcal{C}\to 0$.
\end{relcorrection}

\subsection{General case: Israel--Stewart viscosity}
\label{sec:rel-10-98-general}

Returning to the general viscous case, the solution~\eqref{eq:rel-10-261R}
must satisfy three boundary conditions at $r = R$:

\begin{enumerate}
\item \textbf{Kinematic condition:} the radial velocity matches the
  surface displacement,
  \begin{equation}\label{eq:rel-10-264R}
    l(l+1)\!\left[
      A\,\frac{J_{l+\frac{1}{2}}(z_{\mathrm{rel}})}{R^{1/2}}
      + \frac{\epsilon_{0}}{\nu_{\mathrm{eff}}}\,
        \Pi_{0}^{\mathrm{rel}}\,R^{l-1}
    \right]
    = -\epsilon_{0}\,\sigma,
  \end{equation}
  where $z_{\mathrm{rel}} = q_{\mathrm{rel}}\,R$.

\item \textbf{Tangential stress:} the Israel--Stewart shear stress
  $\delta\pi^{r\theta}$ vanishes at the free surface,
  \begin{equation}\label{eq:rel-10-266R}
    \frac{l(l+1)}{r}\,U - \frac{2}{r}\,\frac{dU}{dr}
    + \frac{d^{2}U}{dr^{2}} = 0
    \quad\text{at}\;r = R.
  \end{equation}
  This has the same form as classically because the Israel--Stewart
  prefactor $(1 - \taupi\,\sigma)^{-1}$ cancels at a free boundary
  where $\delta\pi^{r\theta} = 0$.

\item \textbf{Normal stress:} the total $(r,r)$-component of the
  stress-energy tensor (including the Israel--Stewart correction)
  vanishes on the deformed boundary,
  \begin{equation}\label{eq:rel-10-272R}
    -\frac{dp_{0}}{dr}\bigg|_{R}\,\epsilon\,Y_{l}^{m}
    + \left(\frac{\delta p}{\enthalpy/c^{2}}\right)_{\!R}
    - \frac{2\,\nu_{\mathrm{eff}}\,(\enthalpy/c^{2})}{(\enthalpy/c^{2})}
      \left(\frac{\partial u_{r}}{\partial r}\right)_{\!R}
    = 0.
  \end{equation}
\end{enumerate}

\subsection{Characteristic equation}\label{sec:rel-10-98-chareq}

Eliminating $A$ and $\Pi_{0}^{\mathrm{rel}}$ between the three
boundary conditions and defining
\begin{equation}\label{eq:rel-10-270R}
  Q_{l}(z_{\mathrm{rel}})
  = \frac{J_{l+\frac{1}{2}}(z_{\mathrm{rel}})}
         {J'_{l+\frac{1}{2}}(z_{\mathrm{rel}})},
\end{equation}
we arrive at the relativistic characteristic equation
\begin{equation}\label{eq:rel-10-280R}
  \boxed{
  \frac{\sigma}{\sigma_{Kl,\mathrm{rel}}^{2}}
  = \frac{2(l^{2}-1)}{z_{\mathrm{rel}}^{2}} - 1
    + \frac{2\,l(l-1)}{z_{\mathrm{rel}}^{4}}
      \left[1 - \frac{(l+1)\,Q_{l+\frac{1}{2}}(z_{\mathrm{rel}})}
                      {z_{\mathrm{rel}}^{2}
                       - Q_{l+\frac{1}{2}}(z_{\mathrm{rel}})}
      \right],
  }
\end{equation}
where
\begin{equation}\label{eq:rel-10-zrel}
  z_{\mathrm{rel}}
  = R\,\sqrt{\frac{\sigma\,(1 - \taupi\,\sigma)}{\nu}}.
\end{equation}

Equation~\eqref{eq:rel-10-280R} has the \emph{same functional form}
as the classical equation~(280) of \S\,98, but with two crucial
differences:
\begin{enumerate}
\item The Kelvin frequency is $\sigma_{Kl,\mathrm{rel}}$ rather
  than $\sigma_{Kl}$, incorporating the gravitational relativistic
  correction $\beta_{l}\,\mathcal{C}$.
\item The argument $z_{\mathrm{rel}}$ of the Bessel functions depends
  on $\sigma$ through the Israel--Stewart factor
  $(1 - \taupi\,\sigma)$, which raises the effective order of the
  characteristic equation and introduces additional (transient) modes.
\end{enumerate}

In the limit $\mathcal{C}\to 0$ and $\taupi\to 0$, we recover
$\sigma_{Kl,\mathrm{rel}}\to\sigma_{Kl}$ and
$z_{\mathrm{rel}}\to z = R\sqrt{\sigma/\nu}$, thus reproducing
equation~(280) of \S\,98 identically.

\subsection{Decay of Kelvin modes --- relativistic damping rates}
\label{sec:rel-10-98-decay}

\subsubsection*{Small viscosity: $\nu\to 0$ (Israel--Stewart)}

In the limit of small viscosity $\nu\to 0$ with $\taupi$ fixed,
$|z_{\mathrm{rel}}|\to\infty$.  Using the asymptotic expansion of
$Q_{l+1/2}$ for large argument, the characteristic
equation~\eqref{eq:rel-10-280R} yields
\begin{equation}\label{eq:rel-10-290R}
  \sigma
  = \pm\,i\,\sigma_{Kl,\mathrm{rel}}
    + (l-1)(2l+1)\,\frac{\nu}{R^{2}}
      \left[1 + \taupi\,\sigma_{Kl,\mathrm{rel}}^{2}\,
                \frac{\nu}{R^{2}}
            + \order{\taupi^{2}\sigma_{Kl}^{4}\nu^{2}/R^{4}}
      \right]
  \quad (\nu\to 0).
\end{equation}
The mean life of the Kelvin $l$-mode is therefore
\begin{equation}\label{eq:rel-10-291R}
  \boxed{
  \tau_{l}^{\mathrm{rel}}
  = \frac{R^{2}}{(l-1)(2l+1)\,\nu}
    \left[1 - \taupi\,\sigma_{Kl,\mathrm{rel}}^{2}\,
              \frac{\nu}{R^{2}}
          + \order{\taupi^{2}}
    \right].
  }
\end{equation}
The Israel--Stewart correction \emph{shortens} the decay time
compared with the Navier--Stokes prediction: the causal delay
in establishing viscous stresses makes the damping marginally more
efficient at second order.

\begin{relcorrection}
Classically (equation~(291) of \S\,98),
$\tau_{l} = R^{2}/[(l-1)(2l+1)\nu]$.
The leading relativistic correction to the damping rate has two
parts: (a) the shift in $\sigma_{Kl}$ by a factor
$(1 + \beta_{l}\,\mathcal{C}/2)$, and (b) the Israel--Stewart
correction proportional to $\taupi$.  Both vanish in the Newtonian,
instantaneous-viscosity limit.
\end{relcorrection}

\subsubsection*{Large viscosity: $\nu\to\infty$}

For $\nu\to\infty$ (the ``creeping'' regime),
$z_{\mathrm{rel}}\to 0$ and
\begin{equation}\label{eq:rel-10-287R}
  \sigma \to
  \sigma_{Kl,\mathrm{rel}}^{2}\,
  \frac{(2l+1)\,R^{2}}{2(l-1)(2l^{2}+4l+3)\,\nu}
  \quad (\nu\to\infty),
\end{equation}
which has the same structure as the classical
equation~(287) of \S\,98 but with
$\sigma_{Kl}\to\sigma_{Kl,\mathrm{rel}}$.  The Israel--Stewart
relaxation time drops out of the leading term in this limit because
$\taupi\,\sigma \to 0$ when $\sigma\to 0$.

\subsubsection*{Higher-order aperiodic modes}

As in the classical analysis, additional aperiodic modes arise
from the poles of $Q_{l+1/2}(z_{\mathrm{rel}})$ at the zeros
$x_{l+1/2,j}$ of $J'_{l+1/2}$.  To leading order these are
characterised by
\begin{equation}\label{eq:rel-10-288R}
  \sigma \sim
  \frac{x_{l+\frac{1}{2},j}^{2}\,\nu}{R^{2}\,(1 + \taupi\,x_{l+\frac{1}{2},j}^{2}\,\nu/R^{2})}
  \quad (j \geqslant 2),
\end{equation}
which reduces to equation~(288) of \S\,98 when $\taupi\to 0$.
The Israel--Stewart relaxation acts as a saturation mechanism:
for large $j$, the damping constant approaches $1/\taupi$
rather than growing without bound, consistent with causality.

\begin{causalitycheck}
In the Navier--Stokes theory ($\taupi = 0$), the higher-order
damping constants $\sigma_{j} \sim x_{j}^{2}\,\nu/R^{2}$ grow
without bound, implying arbitrarily fast signal propagation ---
a well-known pathology.  The Israel--Stewart
result~\eqref{eq:rel-10-288R} caps all damping rates at
$\sigma_{\max} = 1/\taupi$, ensuring that the corresponding
propagation speed $v_{\mathrm{shear}} = R\,\sqrt{\sigma/\nu_{\mathrm{eff}}}$
never exceeds $c$.
\end{causalitycheck}

%% ================================================================

\section{Oscillations of a relativistic viscous liquid drop
(\S99 -- relativistic)}\label{sec:rel-10-99}

We now consider the relativistic extension of \S\,99: the oscillations
of a viscous liquid drop held together by surface tension rather than
self-gravity.  As Chandrasekhar and Reid showed classically, when the
decay rate $\sigma$ is expressed in units of the inviscid eigenfrequency
$\sigma_{Kl}$, the characteristic equations for the globe and the drop
become identical.  We show that this remarkable correspondence persists
in the relativistic theory, provided the surface tension is promoted to
a relativistic surface energy.

\subsection{Relativistic surface energy and the Israel junction
conditions}\label{sec:rel-10-99-surface}

In special relativity, a surface with tension $T$ (energy per unit area)
contributes to the energy-momentum tensor via the Israel junction
conditions.  For a spherical drop with surface deformation
$r = R + \epsilon\,Y_{l}^{m}$, the sum of principal curvatures is
\begin{equation}\label{eq:rel-10-298R}
  \frac{1}{R_{1}} + \frac{1}{R_{2}}
  = \frac{2}{R}
    + \frac{\epsilon}{R^{2}}\,(l-1)(l+2)\,Y_{l}^{m},
\end{equation}
exactly as in the classical case (equation~(298) of \S\,99).
However, the boundary condition on the stress-energy tensor now reads
\begin{equation}\label{eq:rel-10-302R}
  \left[T^{\mu\nu}\right]_{\Sigma}\,n_{\nu}
  = -S^{\mu\nu}\,K_{\nu}^{\phantom{\nu}\alpha}\,n_{\alpha},
\end{equation}
where $S^{\mu\nu}$ is the surface energy-momentum tensor,
$K_{\nu}^{\phantom{\nu}\alpha}$ is the extrinsic curvature, and
$n_{\mu}$ is the surface normal.  For a perfect-fluid surface layer
with surface energy density $\edensity_{S}$ and surface tension $T$,
the surface energy-momentum tensor is
\begin{equation}\label{eq:rel-10-Smunu}
  S^{\mu\nu}
  = (\edensity_{S} + T)\,\bar{u}^{\mu}\,\bar{u}^{\nu}/c^{2}
    + T\,h^{\mu\nu},
\end{equation}
where $h^{\mu\nu}$ is the induced metric on the surface and
$\bar{u}^{\mu}$ is the surface 4-velocity.

For a \emph{static} equilibrium surface, the Israel condition
reduces to the relativistic Young--Laplace equation
\begin{equation}\label{eq:rel-10-YL-rel}
  \Delta p = T\!\left(\frac{1}{R_{1}} + \frac{1}{R_{2}}\right)
  + \frac{\edensity_{S}}{R},
\end{equation}
where the last term is the relativistic correction arising from the
surface energy density.

\subsection{Inviscid case}\label{sec:rel-10-99-inviscid}

Setting $\nu = 0$ and repeating the boundary analysis with the
relativistic Young--Laplace equation, the inviscid eigenfrequency
becomes
\begin{equation}\label{eq:rel-10-301R}
  \boxed{
  \sigma_{Kl,\mathrm{drop}}^{2}
  = l(l-1)(l+2)\,\frac{T}{\enthalpy\,R^{3}/c^{2}}
    \left[1 + \frac{(l-1)\,\edensity_{S}}{T\,R}
          + \order{v^{2}/c^{2}}\right].
  }
\end{equation}
Compared with the classical result
$\sigma_{Kl}^{2} = l(l-1)(l+2)\,T/(\rho\,R^{3})$
(equation~(301) of \S\,99), two modifications appear:
\begin{enumerate}
\item The inertial density is
  $\enthalpy/c^{2} = (\edensity + p)/c^{2}$ rather than $\rho$.
\item The surface energy density $\edensity_{S}$ provides an additional
  restoring contribution.
\end{enumerate}
Both corrections vanish in the non-relativistic limit
$\edensity\to\rho\,c^{2}$, $p \ll \rho\,c^{2}$,
$\edensity_{S}\ll T$.

\subsection{General case with Israel--Stewart viscosity}
\label{sec:rel-10-99-general}

For the viscous drop, the perturbation
equations~\eqref{eq:rel-10-236R}--\eqref{eq:rel-10-237R} remain
valid, but with $\delta V = 0$ (no self-gravity) and
$\varpi_{\mathrm{rel}} = \delta p/(\enthalpy/c^{2})$.
The velocity field and the Israel--Stewart shear stress are
identical to those of the globe problem.  The only change is in
the normal-stress boundary condition, where the gravitational
restoring force is replaced by the relativistic surface tension
condition~\eqref{eq:rel-10-302R}.

Following the same algebra as in
\S\,\ref{sec:rel-10-98-chareq}, the boundary conditions yield
\begin{equation}\label{eq:rel-10-304R}
  \left\{-\sigma_{Kl,\mathrm{drop}}^{2}
         + l(l+1)\,\Pi_{0}^{\mathrm{drop}}\,R^{l-1}
  \right\}\epsilon\,Y_{l}^{m}
  - \frac{2\,\nu_{\mathrm{eff}}}{R}
    \left(\frac{\partial u_{r}}{\partial r}\right)_{\!R}
  = 0,
\end{equation}
which has \emph{precisely the same form} as the globe
condition~\eqref{eq:rel-10-272R} with the replacement
$\sigma_{Kl,\mathrm{rel}}\to\sigma_{Kl,\mathrm{drop}}$.

\subsection{Identity of the characteristic equations}

Since equation~\eqref{eq:rel-10-304R} and the kinematic and
tangential boundary conditions are structurally identical for the
globe and the drop, the elimination of $A$ and $\Pi_{0}$ yields
\emph{the same characteristic equation}~\eqref{eq:rel-10-280R}
with $\sigma_{Kl,\mathrm{rel}}$ replaced by
$\sigma_{Kl,\mathrm{drop}}$.

\begin{relcorrection}
The Chandrasekhar--Reid identity --- that the globe and drop problems
share the same characteristic equation when $\sigma$ is measured in
units of the appropriate Kelvin frequency --- survives intact in
the relativistic, Israel--Stewart theory.  The forces maintaining
equilibrium (gravity or surface tension) enter only through
$\sigma_{Kl}$; the manner of viscous decay is universal.  The
Israel--Stewart relaxation modifies the functional form of the
characteristic equation (through $z_{\mathrm{rel}}$) but does so
identically for both problems.
\end{relcorrection}

All results for damping rates, creeping modes, and higher-order
aperiodic modes derived in
\S\,\ref{sec:rel-10-98-decay} carry over to the drop upon the
substitution
\begin{equation}\label{eq:rel-10-drop-sub}
  \sigma_{Kl,\mathrm{rel}}^{2}
  \;\longrightarrow\;
  \sigma_{Kl,\mathrm{drop}}^{2}
  = l(l-1)(l+2)\,\frac{T}{\enthalpy\,R^{3}/c^{2}}
    \left[1 + \frac{(l-1)\,\edensity_{S}}{T\,R}\right].
\end{equation}

For a relativistic drop (e.g.\ a quark-gluon plasma droplet) of
radius $R$ with surface tension $T$ and surface energy
$\edensity_{S} \sim T$, the Israel condition correction
$(l-1)\edensity_{S}/(T\,R) \sim (l-1)/R$ becomes significant
when $R$ is comparable to the microscopic scale $\edensity_{S}/T$.

%% ================================================================

\section*{Bibliographical Notes}

\noindent\S\,98 (relativistic).
The classical oscillations of a viscous liquid globe were first
treated by Lamb~\cite{Lamb1881} and developed in the form presented
here by Chandrasekhar~\cite{Chandrasekhar1959globe}.  The relativistic
generalisation of stellar oscillation theory, in the inviscid case,
was initiated by Thorne and Campolattaro:

\begin{enumerate}
\item \textsc{K.\,S.\ Thorne} and \textsc{A.\ Campolattaro},
  `Non-radial pulsation of general-relativistic stellar models.\ I.
  Analytic analysis for $l \geqslant 2$',
  \textit{Astrophys.\ J.}\ \textbf{149}, 591--611 (1967).
\end{enumerate}

\noindent
The Israel--Stewart framework for causal relativistic dissipation is
due to:

\begin{enumerate}
\setcounter{enumi}{1}
\item \textsc{W.\ Israel} and \textsc{J.\,M.\ Stewart},
  `Transient relativistic thermodynamics and kinetic theory',
  \textit{Ann.\ Phys.\ (N.Y.)}\ \textbf{118}, 341--72 (1979).
\end{enumerate}

\noindent
Viscous damping of stellar oscillations in general relativity, using
the Israel--Stewart formalism, has been studied by:

\begin{enumerate}
\setcounter{enumi}{2}
\item \textsc{L.\ Lindblom} and \textsc{S.\,L.\ Detweiler},
  `The effect of viscosity on oscillations of a uniformly
  rotating star',
  \textit{Astrophys.\ J.}\ \textbf{292}, 12--15 (1985).

\item \textsc{N.\ Andersson} and \textsc{G.\,L.\ Comer},
  `Relativistic fluid dynamics: physics for many different scales',
  \textit{Living Rev.\ Relativ.}\ \textbf{10}, 1 (2007).
\end{enumerate}

\noindent
The Tolman--Oppenheimer--Volkoff equation for relativistic stellar
equilibrium was derived in:

\begin{enumerate}
\setcounter{enumi}{4}
\item \textsc{J.\,R.\ Oppenheimer} and \textsc{G.\,M.\ Volkoff},
  `On massive neutron cores',
  \textit{Phys.\ Rev.}\ \textbf{55}, 374--81 (1939).
\end{enumerate}

\medskip
\noindent\S\,99 (relativistic).
The classical identity between the globe and drop characteristic
equations was established by Reid~\cite{Reid1960drop}.  The
Israel junction conditions for matching spacetimes across a
surface layer are due to:

\begin{enumerate}
\setcounter{enumi}{5}
\item \textsc{W.\ Israel},
  `Singular hypersurfaces and thin shells in general relativity',
  \textit{Nuovo Cim.}\ B\,\textbf{44}, 1--14 (1966);
  erratum: ibid.\ \textbf{48}, 463 (1967).
\end{enumerate}

\noindent
The relativistic surface tension of quark-gluon plasma droplets
is discussed in:

\begin{enumerate}
\setcounter{enumi}{6}
\item \textsc{M.\,G.\ Alford}, \textsc{K.\ Rajagopal},
  and \textsc{F.\ Wilczek},
  `Color-flavor locking and chiral symmetry breaking in
  high density QCD',
  \textit{Nucl.\ Phys.}\ B\,\textbf{537}, 443--58 (1999).

\item \textsc{P.\ Romatschke},
  `New developments in relativistic fluid dynamics',
  \textit{Int.\ J.\ Mod.\ Phys.}\ E\,\textbf{19}, 1--53 (2010).
\end{enumerate}
